\documentclass[12pt,a4paper]{article}
\usepackage[utf8]{inputenc}
\usepackage[T1]{fontenc}
\usepackage{amsmath}
\usepackage[left=2.00cm, right=2.00cm, top=2.00cm, bottom=2.00cm]{geometry}
\usepackage{amssymb}
\usepackage[italian]{babel}

\title{Soluzioni Tutorato 3\\Complementi di Matematica per le Scienze Chimiche\\ A.A. 2021/2022}
\author{prof. Emanuele Bottazzi}
\date{}

\usepackage{amsthm}
\theoremstyle{definition}
\newtheorem{exe}{Esercizio}
\newtheorem*{sol}{Soluzione}

\begin{document}
\maketitle

\begin{exe}
    Determina l'insieme aperto $A \subseteq \mathbb{R}^3$ dove le funzioni sono di classe $C^1$.
\end{exe}
\begin{sol}
    \begin{enumerate}
        \item $A = \mathbb{R}^3$;
        \item $A = \{(x,y,z) \in \mathbb{R}^3 : z \neq k\pi, k \in \mathbb{Z}\}$;
        \item $A = \{(x,y,z) \in \mathbb{R}^3 : x \neq 0\}$;
        \item $A = \{(x,y,z) \in \mathbb{R}^3 : y \neq -z^2\}$.
    \end{enumerate}
\end{sol}

\begin{exe}
    Calcola $\frac{\partial f(0,0,1)}{\partial v}$ e $\frac{\partial f(0,0,0)}{\partial v}$ con $f(x,y,z) = \cos(xy)(z-1)^3$.
\end{exe}
\begin{sol}
    $f$ è di classe $C^1$ su tutto $\mathbb{R}^3$, quindi possiamo applicare la formula del gradiente.
    $\frac{\partial f}{\partial x} = -y \sin(xy)(z-1)^3$, $\frac{\partial f}{\partial y} = -x \sin(xy)(z-1)^3$, $\frac{\partial f}{\partial z} = 3 \cos(xy)(z-1)^2$.
    $\nabla f(0,0,1) = (0,0,0)$, quindi $D_v f(0,0,1) = 0$ per ogni versore $v$.
    $\nabla f(0,0,0) = (0,0,3)$, quindi $D_v f(0,0,0) = 3v_z$.
\end{sol}

\begin{exe}
    Insieme aperto $A \subseteq \mathbb{R}^2$ dove le funzioni sono $C^2$.
\end{exe}
\begin{sol}
    \begin{enumerate}
        \item $A = \mathbb{R}^2$;
        \item $A = \{(x,y) \in \mathbb{R}^2 : x \neq 0 \wedge y \neq 0\}$;
        \item $A = \{(x,y) \in \mathbb{R}^2 : x \neq 0\}$;
        \item $A = \{(x,y) \in \mathbb{R}^2 : y \neq x^2+1\}$.
    \end{enumerate}
\end{sol}

\begin{exe}
    Esercizio su $f(x,y) = x \sqrt[3]{y}$ e $g_i$.
\end{exe}
\begin{sol}
    Per $i=3$, $g_3(f(x,y)) = x^3 y$, che è polinomiale e quindi differenziabile ovunque.
\end{sol}

\begin{center}
    prof. Emanuele Bottazzi\\
    emanuele.bottazzi@unipv.it
\end{center}

\end{document}
